\section{METR Time Horizon: usar el tiempo de finalización humano para calibrar la dificultad de las tareas}

\subsection{Definición del 50\% time horizon}

\videofigure{figures/time-horizon-definition.jpg}{Time horizon se define como la duración de finalización humana correspondiente a las tareas que el modelo puede completar con una tasa de éxito dada.}{00:04:26--00:06:08}

Denotemos como $t$ el tiempo que un humano competente necesita para completar la tarea, y como $P_\theta(\text{success}\mid t)$ la probabilidad de éxito del modelo. METR usa una logistic curve para ajustar la variación de la tasa de éxito con la longitud de la tarea:

$$
P_\theta(\text{success}\mid t)=
\frac{1}{1+\exp\left[\beta\left(\log t-\log H_{50}\right)\right]}.
$$

\begin{itemize}
    \item $t$: el tiempo de finalización de un humano competente;
    \item $H_{50}$: la duración de tarea correspondiente al punto en que el modelo alcanza una tasa de éxito del 50\%;
    \item $\beta$: la pendiente con la que la tasa de éxito disminuye conforme aumenta la longitud de la tarea;
    \item $P_\theta$: la probabilidad de éxito empírica bajo una configuración de agent dada.
\end{itemize}

Time horizon no es cuánto tiempo estuvo funcionando el modelo, sino \textbf{cuánto tiempo le tomaría a un humano completar la tarea que el modelo puede completar}. Aquí el tiempo humano funciona como una regla empírica de la complejidad de la tarea.

\begin{importantbox}{El tiempo humano es un proxy de dificultad, no de valor}
Una tarea mecánica de organización que toma tres horas puede tener muy poco valor, mientras que un juicio de alto riesgo de diez minutos puede tener mucho valor. Time horizon mide la extensión de la ejecución autónoma y no puede equipararse directamente al impacto económico.
\end{importantbox}

\subsection{El conjunto de tareas y el proceso de medición}

\videofigure{figures/metr-task-suite.jpg}{METR combina operaciones atómicas del orden de segundos, tareas de software e investigación que van de minutos a horas, y ML research tasks más largas.}{00:06:08--00:07:08}

\videofigure{figures/metr-method.jpg}{El proceso de evaluación incluye construir tareas de múltiples escalas, registrar la línea base humana, ejecutar el agent, estimar la tasa de éxito y ajustar el time horizon.}{00:07:08--00:11:16}

Las tareas cubren ingeniería de software, ciberseguridad, razonamiento general e investigación en aprendizaje automático. Para cada tarea, primero se hace que un profesional competente la complete de manera independiente y se registra el tiempo; después se hace que el agent la intente varias veces en un entorno controlado, y el éxito se determina mediante pruebas automáticas o reglas de evaluación manual.

\begin{knowledgebox}{Por qué se ajusta con log-time}
La duración de las tareas abarca varios órdenes de magnitud, desde segundos hasta minutos y horas. Ajustar sobre $\log t$ permite que tareas de distintas escalas se distribuyan de manera más uniforme en la curva, y también hace que «duplicar la capacidad» se manifieste como una tendencia aproximadamente lineal.
\end{knowledgebox}

\videofigure{figures/success-vs-human-time.jpg}{La tasa de éxito del modelo muestra una correlación negativa clara con el tiempo de finalización humano, lo que indica que la duración humana explica una parte considerable de la dificultad para el agent.}{00:12:22--00:14:02}

Aun así, la figura muestra residuales considerables: algunas tareas que toman mucho tiempo pero cuyos pasos son regulares resultan relativamente fáciles para el agent, mientras que algunas tareas cortas que requieren juicio fino u operación de herramientas frágiles resultan, por el contrario, difíciles. Por lo tanto, time horizon es una tendencia agregada, no un predictor para cada tarea individual.

\subsection{Resumen de la sección}

METR usa el tiempo humano para mapear tareas heterogéneas a un eje unificado de dificultad, y a partir de la curva de tasa de éxito lee $H_{50}$ o $H_{80}$. Es adecuado para seguir tendencias de capacidad a largo plazo, pero es necesario conservar el análisis por tipo de tarea y por modo de falla.
